\documentclass[../../main.tex]{subfiles}

\begin{document}

\section{外测度诱导的 Carath\'eodory 测度}\label{section:ch17-caratheodory}

我们现在定义外测度与集合关于外测度的可测性的一般概念, 且证明构造实直线上的 Lebesgue 测度的 Carath\'eodory 策略一般是可行的.

\begin{definition}[可数单调的集函数]\label{definition:ch17-countably-subadditive}
定义在集合 $X$ 的子集族 $\mathcal{S}$ 上的集函数 $\mu:\mathcal{S}\to[0,\infty]$ 称为可数单调的, 只要集合 $E\in\mathcal{S}$ 被 $\mathcal{S}$ 中的集合的可数族 $\{E_{k}\}_{k=1}^{\infty}$ 覆盖, 则
\begin{align*}
\mu(E)\leqslant\sum_{k=1}^{\infty}\mu(E_{k}).
\end{align*}
\end{definition}

如同我们已观察到的, 测度的单调性和可数可加性告诉我们测度是可数单调的. 若可数单调集函数 $\mu:\mathcal{S}\to[0,\infty]$ 具有性质 $\varnothing$ 属于 $\mathcal{S}$ 且 $\mu(\varnothing)=0$, 则 $\mu$ 在以下意义下是有限单调的: 只要 $E\in\mathcal{S}$ 被 $\mathcal{S}$ 中的集合的有限族 $\{E_{k}\}_{k=1}^{n}$ 覆盖, 则
\begin{align*}
\mu(E)\leqslant\sum_{k=1}^{n}\mu(E_{k}).
\end{align*}
为看到这一点, 对 $k>n$ 设 $E_{k}=\varnothing$. 特别地, 这样的集函数 $\mu$ 在以下意义下是单调的: 若 $A$ 和 $B$ 属于 $\mathcal{S}$ 且 $A\subseteq B$, 则 $\mu(A)\leqslant\mu(B)$.

\begin{definition}[外测度]\label{definition:ch17-outer-measure}
集函数 $\mu^*:2^{X}\to[0,\infty]$ 称为外测度, 若 $\mu^*(\varnothing)=0$ 且 $\mu^*$ 是可数单调的.
\end{definition}

受到从实直线上的 Lebesgue 外测度构造 Lebesgue 测度的经验的指引, 我们遵循 Constantine Carath\'eodory 的方法, 定义集合的可测性如下.

\begin{definition}[Carath\'eodory 可测]\label{definition:ch17-caratheodory-measurable}
对于外测度 $\mu^*:2^{X}\to[0,\infty]$, 我们称 $X$ 的子集 $E$ 可测(关于 $\mu^*$), 若对 $X$ 的每个子集 $A$,
\begin{align*}
\mu^*(A)=\mu^*(A\cap E)+\mu^*(A\cap E^{c}).
\end{align*}
\end{definition}

由于 $\mu^*$ 是有限单调的, 为证明 $E\subseteq X$ 可测仅须对所有使得 $\mu^*(A)<\infty$ 的 $A\subseteq X$ 证明
\begin{align*}
\mu^*(A)\geqslant\mu^*(A\cap E)+\mu^*(A\cap E^{c}).
\end{align*}

直接从定义看到, $X$ 的一个集合 $E$ 是可测的当且仅当它在 $X$ 中的补是可测的, 且根据 $\mu^*$ 的单调性, 每个外测度为零的集合是可测的. 在这里以及本节的后面部分, $\mu^*:2^{X}\to[0,\infty]$ 是外测度且可测指的是关于 $\mu^*$ 可测.

\begin{proposition}\label{proposition:ch17-finite-union-measurable}
有限可测集族的并是可测的.
\end{proposition}

\begin{proof}
我们首先证明两个可测集的并是可测的. 令 $E_{1}$ 和 $E_{2}$ 为可测集. 令 $A$ 为 $X$ 的任何子集. 首先利用 $E_{1}$ 的可测性, 接着用 $E_{2}$ 的可测性, 我们有
\begin{align*}
\mu^*(A)=\mu^*(A\cap E_{1})+\mu^*(A\cap E_{1}^{c}).
\end{align*}
再次利用 $E_{2}$ 的可测性, 对集合 $A\cap E_{1}^{c}$ 有
\begin{align*}
\mu^*(A\cap E_{1}^{c})=\mu^*([A\cap E_{1}^{c}]\cap E_{2})+\mu^*([A\cap E_{1}^{c}]\cap E_{2}^{c}).
\end{align*}
现在用集合等式
\begin{align*}
[A\cap E_{1}^{c}]\cap E_{2}^{c}=A\cap[E_{1}\cup E_{2}]^{c}
\end{align*}
与
\begin{align*}
[A\cap E_{1}]\cup[A\cap E_{2}\cap E_{1}^{c}]=A\cap[E_{1}\cup E_{2}]
\end{align*}
以及外测度的有限单调性得到
\begin{align*}
\mu^*(A)&=\mu^*(A\cap E_{1})+\mu^*(A\cap E_{1}^{c})\\
&=\mu^*(A\cap E_{1})+\mu^*([A\cap E_{1}^{c}]\cap E_{2})+\mu^*([A\cap E_{1}^{c}]\cap E_{2}^{c})\\
&=\mu^*(A\cap E_{1})+\mu^*([A\cap E_{1}^{c}]\cap E_{2})+\mu^*(A\cap[E_{1}\cup E_{2}]^{c})\\
&\geqslant\mu^*(A\cap[E_{1}\cup E_{2}])+\mu^*(A\cap[E_{1}\cup E_{2}]^{c}).
\end{align*}
因此 $E_{1}\cup E_{2}$ 是可测的. 现在令 $\{E_{k}\}_{k=1}^{n}$ 为任何可测集的有限族. 对一般的 $n$, 我们用归纳法证明并集 $\bigcup_{k=1}^{n}E_{k}$ 的可测性. 对于 $n=1$ 这是平凡的. 假定当 $n-1$ 时它成立. 于是, 由于
\begin{align*}
\bigcup_{k=1}^{n}E_{k}=\left[\bigcup_{k=1}^{n-1}E_{k}\right]\cup E_{n},
\end{align*}
且两个可测集的并是可测的, 集合 $\bigcup_{k=1}^{n}E_{k}$ 是可测的.
\end{proof}

\begin{proposition}\label{proposition:ch17-finite-additivity}
令 $A\subseteq X$ 而 $\{E_{k}\}_{k=1}^{n}$ 是可测集的有限不交族. 则
\begin{align*}
\mu^*\left(A\cap\left[\bigcup_{k=1}^{n}E_{k}\right]\right)=\sum_{k=1}^{n}\mu^*(A\cap E_{k}).
\end{align*}
特别地, $\mu^*$ 对可测集族的限制是有限可加的.
\end{proposition}

\begin{proof}
通过对 $n$ 归纳进行证明. $n=1$ 时命题显然成立, 我们假设命题对 $n-1$ 成立. 由于族 $\{E_{k}\}_{k=1}^{n}$ 是不交的,
\begin{align*}
A\cap\left[\bigcup_{k=1}^{n}E_{k}\right]\cap E_{n}=A\cap E_{n}
\end{align*}
且
\begin{align*}
A\cap\left[\bigcup_{k=1}^{n}E_{k}\right]\cap E_{n}^{c}=A\cap\left[\bigcup_{k=1}^{n-1}E_{k}\right].
\end{align*}
因此根据 $E_{n}$ 的可测性以及归纳假设, 我们有
\begin{align*}
\mu^*\left(A\cap\bigcup_{k=1}^{n}E_{k}\right)&=\mu^*(A\cap E_{n})+\mu^*\left(A\cap\bigcup_{k=1}^{n-1}E_{k}\right)\\
&=\mu^*(A\cap E_{n})+\sum_{k=1}^{n-1}\mu^*(A\cap E_{k})\\
&=\sum_{k=1}^{n}\mu^*(A\cap E_{k}).
\end{align*}
\end{proof}

\begin{proposition}\label{proposition:ch17-countable-union-measurable}
可测集的可数族的并是可测的.
\end{proposition}

\begin{proof}
令 $E=\bigcup_{k=1}^{\infty}E_{k}$, 其中每个 $E_{k}$ 是可测的. 由于可测集在 $X$ 中的补是可测的, 根据命题\cref{proposition:ch17-finite-union-measurable}, 可测集的有限族的并是可测的, 通过用 $E_{k}\sim\bigcup_{i=1}^{k-1}E_{i}$ 代替每个 $E_{k}$, 我们可以假定 $\{E_{k}\}_{k=1}^{\infty}$ 是不交的. 令 $A$ 为 $X$ 的任何子集. 固定指标 $n$. 定义 $F_{n}=\bigcup_{k=1}^{n}E_{k}$. 由于 $F_{n}$ 可测且 $F_{n}^{c}\supseteq E^{c}$, 我们有
\begin{align*}
\mu^*(A)=\mu^*(A\cap F_{n})+\mu^*(A\cap F_{n}^{c})\geqslant\mu^*(A\cap F_{n})+\mu^*(A\cap E^{c}).
\end{align*}
根据命题\cref{proposition:ch17-finite-additivity},
\begin{align*}
\mu^*(A\cap F_{n})=\sum_{k=1}^{n}\mu^*(A\cap E_{k}).
\end{align*}
因此,
\begin{align*}
\mu^*(A)\geqslant\sum_{k=1}^{n}\mu^*(A\cap E_{k})+\mu^*(A\cap E^{c}).
\end{align*}
该不等式的左边与 $n$ 无关, 因此,
\begin{align*}
\mu^*(A)\geqslant\sum_{k=1}^{\infty}\mu^*(A\cap E_{k})+\mu^*(A\cap E^{c}).
\end{align*}
根据外测度的可数单调性, 我们推出
\begin{align*}
\mu^*(A)\geqslant\mu^*(A\cap E)+\mu^*(A\cap E^{c}).
\end{align*}
因此 $E$ 是可测的.
\end{proof}

\begin{theorem}\label{theorem:ch17-measure-space}
令 $\mu^*$ 为 $2^{X}$ 上的外测度. 则关于 $\mu^*$ 可测的集族 $\mathcal{M}$ 是一个 $\sigma$ 代数. 若 $\overline{\mu}$ 是 $\mu^*$ 在 $\mathcal{M}$ 上的限制, 则 $(X,\mathcal{M},\overline{\mu})$ 是一个完备的测度空间.
\end{theorem}

\begin{proof}
我们已经观察到 $X$ 的可测子集在 $X$ 中的补也是可测的. 根据命题\cref{proposition:ch17-countable-union-measurable}, 可测集的可数族的并是可测的. 因此 $\mathcal{M}$ 是一个 $\sigma$ 代数. 根据外测度的定义, $\mu^*(\varnothing)=0$, 因此 $\varnothing$ 是可测的且 $\overline{\mu}(\varnothing)=0$. 为证明 $\overline{\mu}$ 是 $\mathcal{M}$ 上的测度, 还需要证明它是可数可加的. 由于 $\mu^*$ 是可数单调的且集函数 $\mu^*$ 是 $\overline{\mu}$ 的延拓, 集函数 $\overline{\mu}$ 是可数单调的. 因此我们仅须证明: 若 $\{E_{k}\}_{k=1}^{\infty}$ 是可测集的不交族, 则
\begin{align}
\mu^*\left(\bigcup_{k=1}^{\infty}E_{k}\right)\geqslant\sum_{k=1}^{\infty}\mu^*(E_{k}).\label{align:ch17-countable-add-ineq}
\end{align}
然而, $\mu^*$ 是单调的, 在命题\cref{proposition:ch17-finite-additivity}中取 $A=X$, 我们看到在可测集的有限不交族上 $\mu^*$ 是可加的. 因此, 对每个 $n$,
\begin{align*}
\mu^*\left(\bigcup_{k=1}^{\infty}E_{k}\right)\geqslant\mu^*\left(\bigcup_{k=1}^{n}E_{k}\right)=\sum_{k=1}^{n}\mu^*(E_{k}).
\end{align*}
不等式的左边与 $n$ 无关, 因此\eqref{align:ch17-countable-add-ineq}成立.
\end{proof}

\end{document}